\documentclass[11pt,a4paper]{article}

% --- Paquetes básicos ---
\usepackage[utf8]{inputenc}
\usepackage[T1]{fontenc}
\usepackage[spanish]{babel}
\usepackage[a4paper,margin=2.5cm]{geometry}
\usepackage{xcolor}
\usepackage{booktabs}
\usepackage{longtable}
\usepackage{array}
\usepackage{hyperref}
\usepackage{enumitem}
\usepackage{tabularx}
\usepackage{graphicx}
\usepackage{amsmath}

% --- Hipervínculos ---
\hypersetup{
  colorlinks=true,
  linkcolor=blue!60!black,
  urlcolor=blue!60!black,
  citecolor=blue!60!black,
  pdftitle={recap-agent: agente conversacional serverless para Sin Envolturas},
  pdfauthor={Tesis}
}

% --- Estilo para código (usando fancyvrb) ---
\usepackage{fancyvrb}
\definecolor{codebg}{RGB}{248,248,248}
\definecolor{codegray}{RGB}{60,60,60}

\DefineVerbatimEnvironment{codeblock}{Verbatim}{
  fontfamily=tt,
  fontsize=\small,
  frame=single,
  framerule=0pt,
  framesep=8pt,
  xleftmargin=1em,
  xrightmargin=1em,
  formatcom=\color{codegray}
}

\newcommand{\code}[1]{{\ttfamily\color{codegray}#1}}

% Título y autor
\title{\textbf{recap-agent}\\[0.4em]
       \large Runtime conversacional serverless para Sin Envolturas:\\
       arquitectura, decisiones y características clave}
\author{Documentación derivada de \texttt{docs/implementation-log.md} y del código fuente}
\date{Generado el 9 de junio de 2026}

\begin{document}

\maketitle

\begin{abstract}
\noindent
Este documento describe \texttt{recap-agent}, un agente conversacional serverless
desplegado sobre AWS Lambda y modelado en torno a un grafo de estados que
refleja el \emph{Flujo de estados} definido para el producto. La arquitectura
combina el SDK de OpenAI Agents (un agente conversacional más un extractor
estructurado), persistencia por turno en DynamoDB, una puerta de enlace real
contra el marketplace de Sin Envolturas y una búsqueda híbrida (API +
vector store) para proveedores. El runtime es agnóstico al canal: el
comportamiento de WhatsApp, webchat y la terminal CLI se resuelve en una capa
de adaptadores delgada, mientras que el núcleo de la aplicación sólo conoce
\texttt{channel} y \texttt{user\_id} como datos de entrada. El presente
artículo resume la arquitectura del sistema, las decisiones de diseño más
relevantes, las características clave implementadas y la evolución que
recoge el log de implementación mantenido desde el inicio del proyecto.
\end{abstract}

\tableofcontents
\bigskip

% ============================================================
\section{Introducción y motivación}
% ============================================================

\texttt{recap-agent} nace como un \emph{runtime conversacional serverless} cuyo
objetivo es asistir a usuarios de \textbf{Sin Envolturas} en la planificación
de eventos: detectar el tipo de evento, elicitar necesidades de proveedores,
buscar y recomendar, gestionar selecciones múltiples, y finalmente cerrar el
plan enviando cotizaciones a los proveedores elegidos. La motivación original
era doble:

\begin{enumerate}[leftmargin=*,itemsep=2pt]
  \item Reemplazar un cliente terminal improvisado por un servicio desplegado
        que pueda ser invocado por canales reales (WhatsApp, webchat), sin
        sacrificar la trazabilidad de la conversación.
  \item Garantizar que las decisiones de flujo no dependan de coincidencias
        de cadenas en el texto del usuario, sino de evidencia estructurada
        producida por un agente extractor y validada por un esquema de
        estados tipado.
\end{enumerate}

El proyecto se escribe íntegramente en TypeScript, usa el runtime
\texttt{nodejs24.x} tanto en local como en Lambda, aplica tipado estricto con
\texttt{any} explícitamente prohibido, y mantiene sus prompts en español
dentro del repositorio como archivos de texto bajo \texttt{prompts/}
asociados a nodos del grafo de estados.

% ============================================================
\section{Arquitectura general}
% ============================================================

\subsection{Vista de capas}

El sistema se organiza en cinco capas, todas ellas desplegadas en una única
plantilla de CloudFormation (\texttt{infra/cloudformation/stack.yaml}):

\begin{itemize}[leftmargin=*,itemsep=2pt]
  \item \textbf{Capa de canal} (\texttt{src/terminal/}): CLI de desarrollo
        que emula WhatsApp; en producción se reemplaza por adaptadores
        específicos de cada canal, sin tocar el núcleo.
  \item \textbf{Entry point} (\texttt{src/lambda/handler.ts}): handler de
        AWS Lambda que invoca el runtime, serializa el turno y persiste
        telemetría. Es la única superficie HTTP del sistema.
  \item \textbf{Runtime de aplicación} (\texttt{src/runtime/}): orquesta
        extracción, búsqueda, render, persistencia y telemetría. Alberga
        la integración con OpenAI Agents, el gateway de Sin Envolturas y
        los renderers por canal.
  \item \textbf{Núcleo de dominio} (\texttt{src/core/}): grafo de
        decisión, plan persistente tipado, suficiencia, decisión por
        turno, prioridades por tipo de evento y trazabilidad.
  \item \textbf{Almacenamiento} (\texttt{src/storage/}): interfaces
        \texttt{PlanStore} y \texttt{PerfStore} con implementaciones
        DynamoDB e in-memory.
\end{itemize}

\begin{figure}[h!]
\centering
\begin{tabularx}{\linewidth}{|>{\raggedright\arraybackslash}p{3.2cm}|>{\raggedright\arraybackslash}X|}
\hline
\textbf{Capa} & \textbf{Responsabilidad} \\
\hline
Canal & Webhook y presentación específica de cada canal (WhatsApp, webchat, terminal). \\
\hline
Lambda & Punto de entrada HTTP, persistencia de telemetría, respuesta síncrona única por turno. \\
\hline
Runtime & Orquestación del turno: extractor estructurado, búsqueda, render, persistencia del plan, telemetría. \\
\hline
Núcleo & Modelo del plan, nodos de decisión, suficiencia, decisión por turno, prioridades por evento. \\
\hline
Almacenamiento & Persistencia por \texttt{(channel, externalUserId)} en DynamoDB (planes y métricas). \\
\hline
\end{tabularx}
\caption{Resumen de capas del runtime.}
\end{figure}

\subsection{Despliegue en AWS}

La plantilla de CloudFormation (\texttt{recap-agent-runtime}) crea:

\begin{itemize}[leftmargin=*,itemsep=2pt]
  \item Una función Lambda sobre \texttt{nodejs24.x}, con 90 segundos de
        timeout y 1024\,MB de memoria, accesible vía \emph{Function URL}
        sin autenticación (la autenticación por canal es responsabilidad
        del adaptador).
  \item Dos tablas DynamoDB en modo \texttt{PAY\_PER\_REQUEST}:
        \texttt{<stack>-plans} (planes por usuario) y
        \texttt{<stack>-perf} (telemetría con GSI \texttt{channel-user-turns}
        y TTL configurable vía \texttt{PERF\_RETENTION\_DAYS}, por defecto 30 días).
  \item Un rol IAM con permisos mínimos sobre DynamoDB y Secrets Manager
        para resolver la API key de OpenAI en tiempo de ejecución.
  \item Un secreto de OpenAI que se sincroniza desde el \texttt{.env} local
        mediante \texttt{scripts/deploy.mjs}.
\end{itemize}

Stacks auxiliares (\texttt{recap-agent-provider-sync-dev},
\texttt{recap-agent-knowledge-sync-dev}) mantienen los vector stores de
proveedores y de la base de conocimiento respectivamente.

% ============================================================
\section{Núcleo de dominio}
% ============================================================

\subsection{Grafo de decisión}

El corazón del sistema es un \textbf{grafo de estados} explícito definido en
\texttt{src/core/decision-nodes.ts}. Cada turno del usuario es procesado en
uno de 26 nodos posibles, entre los que se incluyen:

\begin{itemize}[leftmargin=*,itemsep=1pt]
  \item \textbf{Entrada}: \texttt{contacto\_inicial}, \texttt{deteccion\_intencion},
        \texttt{existe\_plan\_guardado}.
  \item \textbf{Elicitación}: \texttt{entrevista}, \texttt{elicitacion\_necesidades},
        \texttt{minimos\_para\_buscar}, \texttt{aclarar\_pedir\_faltante},
        \texttt{usuario\_responde}.
  \item \textbf{Búsqueda y recomendación}: \texttt{buscar\_proveedores},
        \texttt{busqueda\_exitosa}, \texttt{hay\_resultados},
        \texttt{recomendar}, \texttt{refinar\_criterios}.
  \item \textbf{Selección y refinamiento}: \texttt{usuario\_elige\_proveedor},
        \texttt{anadir\_a\_proveedores\_recomendados},
        \texttt{seguir\_refinando\_guardar\_plan}, \texttt{continua}.
  \item \textbf{Cierre}: \texttt{accion\_final\_exitosa},
        \texttt{necesidad\_cubierta}, \texttt{crear\_lead\_cerrar},
        \texttt{guardar\_seleccion\_reintentar\_luego},
        \texttt{guardar\_cerrar\_temporalmente},
        \texttt{informar\_error\_reintento}, \texttt{reintentar}.
  \item \textbf{Modos informativos}: \texttt{consultar\_faq} y
        \texttt{consultar\_evento\_invitado}, habilitados como nodos de
        primera clase con sus propios prompts y herramientas.
\end{itemize}

Cada nodo tiene asociado un \emph{manifiesto de prompt} en
\texttt{src/runtime/prompt-manifest.ts} que determina qué archivos del
directorio \texttt{prompts/nodes/<nodo>/} se cargan y qué herramientas
puede invocar el agente conversacional. Esta es la materialización del
principio de ``los prompts son contrato'': el runtime no infiere qué
herramientas están permitidas, las lee del manifiesto.

\subsection{Modelo del plan}

El plan es la unidad durable de estado conversacional. Se modela con Zod
en \texttt{src/core/plan.ts} e incluye, entre otros:

\begin{itemize}[leftmargin=*,itemsep=1pt]
  \item Identidad: \texttt{plan\_id}, \texttt{channel},
        \texttt{external\_user\_id}, \texttt{conversation\_id},
        \texttt{lifecycle\_state} (\texttt{active} | \texttt{finished}).
  \item Estado de flujo: \texttt{current\_node}, \texttt{intent},
        \texttt{intent\_confidence}.
  \item Datos del evento: \texttt{event\_type} (canónico:
        \texttt{boda}, \texttt{cumpleanos}, \texttt{corporativo},
        \texttt{baby\_shower}, \texttt{graduacion}, \texttt{bautizo},
        \texttt{aniversario}, \texttt{quinceanos}, \texttt{otro}),
        \texttt{location}, \texttt{budget\_signal}, \texttt{guest\_range}
        y el árbol de \textbf{necesidades de proveedores}
        (\texttt{provider\_needs}).
  \item Datos de contacto (\texttt{contact\_name}, \texttt{contact\_email},
        \texttt{contact\_phone}) y resumen conversacional
        (\texttt{conversation\_summary}, \texttt{last\_user\_goal},
        \texttt{open\_questions}, \texttt{assumptions}).
  \item Listas planas para compatibilidad (\texttt{recommended\_provider\_ids},
        \texttt{selected\_provider\_ids}, \texttt{selected\_provider\_hints}).
\end{itemize}

La función \texttt{normalizeRawPlan} se ejecuta en el \emph{boundary} de
carga para mapear cualquier valor heredado a los enums canónicos,
incluyendo la migración de los antiguos campos singulares
\texttt{selected\_provider\_id} y \texttt{selected\_provider\_hint} a sus
versiones en arreglo.

\subsection{Suficiencia y \emph{TurnDecision}}

\texttt{computeSearchSufficiency} evalúa si un plan está listo para buscar
proveedores: requiere categoría de proveedor, ubicación, y
\texttt{budget\_signal} o \texttt{guest\_range}. Esta comprobación
determinista se complementa con \texttt{computeNeedSearchSufficiencies},
que aplica la misma regla a cada necesidad del plan.

\texttt{TurnDecision} (\texttt{src/core/turn-decision.ts}) es el
\emph{contrato de ruteo} del turno. Tras la extracción estructurada, el
runtime acumula evidencia (\texttt{DecisionEvidence}) y produce un
\texttt{TurnDecision} validado con Zod que contiene:

\begin{itemize}[leftmargin=*,itemsep=1pt]
  \item \texttt{nextNode}: el nodo que ejecutará la lógica de
        aplicación.
  \item \texttt{routeKind}: tipo de ruta
        (\texttt{ask\_event\_context}, \texttt{clarify\_missing\_fields},
        \texttt{single\_need\_search}, \texttt{multi\_need\_search},
        \texttt{present\_existing\_shortlist}, \texttt{apply\_selection},
        \texttt{modify\_plan}, \texttt{faq}, \texttt{invited\_event\_lookup},
        \texttt{close}, \texttt{pause}, \texttt{error}).
  \item \texttt{providerSearchMode}: \texttt{none},
        \texttt{single\_need\_from\_plan}, \texttt{multi\_need\_query\_intents},
        \texttt{existing\_shortlist}.
  \item \texttt{presentationScope}, \texttt{focusNeedCategory},
        \texttt{needsToSearch}, \texttt{needsToPresent},
        \texttt{invariantStatus} y \texttt{invariantViolations}.
\end{itemize}

El principio clave es: \textbf{el modelo produce interpretación
estructurada; el código de aplicación decide las consecuencias}. De este
modo, el ruteo, la persistencia y la trazabilidad son siempre
deterministas, y se pueden cubrir con tests unitarios.

\subsection{Foco por sesión}

\texttt{sessionFocusSchema} define un objeto persistente que se almacena
como ítem compañero en la tabla de planes. Permite que un mismo usuario
en la misma sesión concentre un frente de trabajo (por ejemplo, ante una
consulta ambigua entre varias necesidades) sin que el runtime caiga en un
estado activo obsoleto. Los adaptadores pueden pasar \texttt{session\_id}
en el cuerpo de la Lambda para habilitar este comportamiento.

% ============================================================
\section{Runtime de aplicación}
% ============================================================

\subsection{Pipeline de un turno}

Cada turno entrante pasa por una pipeline en
\texttt{src/runtime/agent-service.ts} que se puede resumir en:

\begin{enumerate}[leftmargin=*,itemsep=1pt]
  \item Carga del plan desde DynamoDB por
        \texttt{(channel, externalUserId)}.
  \item Validación de invariantes: si el plan está \texttt{finished}, se
        emite respuesta determinista en español sin extractor ni
        composición.
  \item Cálculo de \texttt{DecisionEvidence} y \texttt{TurnDecision}.
  \item Extracción estructurada con el \emph{agente extractor}
        (modelo \texttt{OPENAI\_EXTRACTOR\_MODEL}) usando un esquema Zod
        estable.
  \item Fusión de la extracción con el plan: se preservan hechos
        conocidos (rango de invitados, ubicación, evento) cuando la
        extracción devuelve \texttt{unknown} o nulo.
  \item Aplicación de operaciones estructuradas de plan
        (añadir, actualizar, eliminar, diferir, reactivar, seleccionar,
        reemplazar) y selección de proveedor basada en pista del
        extractor (\texttt{selectedProviderHint}) con respaldo
        determinista de alias/ordinales.
  \item Decisión del modo de búsqueda: ninguno,
        \texttt{single\_need\_from\_plan},
        \texttt{multi\_need\_query\_intents} o
        \texttt{existing\_shortlist}.
  \item Búsqueda y enriquecimiento de proveedores vía el gateway.
  \item Reranking determinista con \texttt{providerFitCriteria}
        (criterios estructurados que el extractor debe producir).
  \item Composición de la respuesta con el \emph{agente conversacional}
        (modelo \texttt{OPENAI\_MODEL}) usando un esquema de salida
        estructurado por nodo (\texttt{welcome}, \texttt{recommendation},
        \texttt{multi\_need\_recommendation}, \texttt{close},
        \texttt{contact\_request}, \texttt{generic}).
  \item Render por canal (\texttt{whatsapp}, \texttt{webchat},
        \texttt{terminal\_whatsapp}).
  \item Persistencia del plan y construcción del registro de telemetría
        (\texttt{buildTurnPerfRecord}).
\end{enumerate}

\subsection{Modelo event-plan-first}

El runtime está explícitamente orientado a un \textbf{plan de evento
primero}, no a una búsqueda puntual. El esquema del plan contiene un
arreglo de \texttt{provider\_needs}, cada una con su propia categoría,
preferencias, restricciones, \texttt{recommended\_providers},
\texttt{sub\_query\_results} y selección múltiple. El campo
\texttt{active\_need\_category} proyecta la necesidad activa para
mantener compatibilidades con el modelo y la trazabilidad.

El catálogo canónico de categorías de proveedor
(\texttt{src/core/provider-category.ts}) se mantiene sincronizado con
los nombres oficiales de la API de Sin Envolturas y se usa como enum Zod
(\texttt{providerCategorySchema}). Adicionalmente, se exponen
\emph{buckets} (Catering, Entretenimiento, Belleza, etc.) que se
expanden a varias categorías canónicas para alimentar búsquedas
paralelas.

\subsection{Gateway de Sin Envolturas}

La clase \texttt{SinEnvolturasGateway}
(\texttt{src/runtime/sinenvolturas-gateway.ts}) implementa el contrato
\texttt{ProviderGateway} y se conecta a la API real del marketplace
(\texttt{https://api.sinenvolturas.com/api-web/vendor}). El conjunto de
operaciones validadas cubre, entre otras:

\begin{itemize}[leftmargin=*,itemsep=1pt]
  \item \texttt{list\_categories}, \texttt{get\_category\_by\_slug},
        \texttt{list\_locations}.
  \item Búsqueda: \texttt{search\_providers\_from\_plan},
        \texttt{search\_providers\_by\_keyword},
        \texttt{search\_providers\_by\_category\_location},
        \texttt{search\_providers\_by\_query\_intent},
        \texttt{get\_relevant\_providers}.
  \item Detalle: \texttt{get\_provider\_detail},
        \texttt{get\_provider\_detail\_and\_track\_view},
        \texttt{get\_related\_providers},
        \texttt{list\_provider\_reviews}.
  \item Contexto: \texttt{get\_event\_vendor\_context},
        \texttt{list\_event\_favorite\_providers},
        \texttt{list\_user\_events\_vendor\_context},
        \texttt{lookup\_user\_event\_context}.
  \item Cierre: \texttt{create\_quote\_request},
        \texttt{add\_vendor\_to\_event\_favorites},
        \texttt{create\_provider\_review},
        \texttt{finish\_plan}.
\end{itemize}

El gateway ofrece una estrategia mixta de búsqueda (API + vector) que
ejecuta en paralelo \texttt{GET /filtered} y \texttt{GET /filtered/full},
deduplica por identificador de proveedor y combina campos
(metadata) para maximizar la completitud sin expandir la superficie de
herramientas que ve el modelo.

\subsection{Búsqueda híbrida con vector store}

La clase \texttt{ProviderVectorSearchGateway}
(\texttt{src/runtime/provider-vector-search.ts}) consulta un vector
store de OpenAI dedicado (\texttt{Sin Envolturas Provider Search})
poblado por la pipeline \texttt{provider-sync} con un archivo Markdown
por proveedor. La estrategia de búsqueda se controla con
\texttt{PROVIDER\_SEARCH\_MODE} (\texttt{api}, \texttt{vector},
\texttt{hybrid}). Las claves canónicas de los filtros de OpenAI se
resuelven en runtime para evitar el problema de \emph{case sensitivity}
que en su momento descartó resultados (por ejemplo, ``Catering'').

El reranking final siempre pasa por el selector determinista de
categoría/ubicación, lo que garantiza que la búsqueda vectorial sólo
\emph{recupera candidatos} y la presentación final se queda con los
proveedores correctos para el país y la ciudad del plan. Esta decisión
se tomó tras observar que un plan de fotografía en Lurín
(Lima/Perú) regresaba candidatos de México.

\subsection{Prompts como contrato}

Cada nodo carga tres archivos desde \texttt{prompts/nodes/<nodo>/}:

\begin{itemize}[leftmargin=*,itemsep=1pt]
  \item \texttt{system.txt}: objetivo, restricciones y comportamiento
        de salida del nodo.
  \item \texttt{response\_contract.txt}: contrato de la salida
        estructurada esperada.
  \item \texttt{tool\_policy.txt}: herramientas permitidas en el nodo.
\end{itemize}

El agente extractor consume, además, un \emph{bundle} propio bajo
\texttt{prompts/extractors/} (\texttt{system},
\texttt{field\_definitions}, \texttt{conflict\_resolution},
\texttt{domain\_knowledge}, \texttt{normalization\_rules},
\texttt{examples.md}), mientras que los prompts compartidos viven en
\texttt{prompts/shared/} (estilo, disciplina de flujo, estrategia de
preguntas, anti-patrones, conocimiento de dominio). El bundle se
inyecta con un \texttt{cacheKey} estable para que la caché de
prompts de OpenAI funcione efectivamente.

\subsection{Guardrails nativos}

El runtime aplica \emph{input} e \emph{output guardrails} del SDK de
OpenAI Agents:

\begin{itemize}[leftmargin=*,itemsep=1pt]
  \item Bloqueo de intentos de prompt injection o de revelación de
        instrucciones internas.
  \item Normalización del correo de soporte al canónico
        \texttt{hola@sinenvolturas.com} si el modelo emite una variante
        corrupta o inventada.
\end{itemize}

Ambas reglas tienen tests de regresión dedicados.

% ============================================================
\section{Características clave implementadas}
% ============================================================

A continuación, una selección de las capacidades más representativas del
runtime tal como se documentan en \texttt{docs/implementation-log.md}.

\begin{longtable}{p{4.6cm}p{10.8cm}}
\toprule
\textbf{Característica} & \textbf{Descripción} \\
\midrule
\endhead

Plan de evento primero &
El plan admite múltiples \texttt{provider\_needs} simultáneas; la
búsqueda de un único proveedor se modela como un caso particular de
este esquema. La categoría activa se proyecta y se normaliza con las
prioridades por tipo de evento. \\

Selección múltiple de proveedores &
\texttt{selected\_provider\_ids} y \texttt{selected\_provider\_hints}
son arreglos. La elección se resuelve con pistas estructuradas del
extractor (alias, ordinales, descriptores como ``la de tablas de
queso'') y con respaldo determinista de matching. La
\texttt{finish\_plan} crea una solicitud de cotización por cada
proveedor seleccionado. \\

Elicitación estructurada multi-necesidad &
El nodo \texttt{elicitacion\_necesidades} produce varios
\texttt{provider\_needs} con \emph{sub-queries} y
\texttt{providerQueryIntents} en una sola pasada cuando la solicitud
del usuario es detallada. La compuerta de \emph{detailed elicitation}
decide si se lanza una búsqueda o un menú inicial. \\

Búsqueda híbrida API + vector store &
Combina los endpoints \texttt{/filtered} y \texttt{/filtered/full} con
un vector store OpenAI dedicado. Deduplica por id, expande buckets
(``Catering'' $\to$ \texttt{[Catering, Licores]}) y aplica un selector
determinista de categoría/ubicación. Los candidatos vectoriales
pasan obligatoriamente por el selector antes de llegar al modelo. \\

Enriquecimiento previo a la recomendación &
Antes de persistir y enviar los resultados al agente conversacional,
el runtime consulta el endpoint de detalle para obtener
\textit{promos}, \textit{service highlights}, \textit{terms} y la URL
de la ficha del proveedor. La presentación final cuenta con
diferenciadores concretos en lugar de texto improvisado. \\

Embudo de recomendación de dos etapas &
El runtime conserva y envía al agente de respuesta un top-15
persistente, y el prompt restringe la presentación visible a un
top-5 (\texttt{PRESENTATION\_PROVIDER\_LIMIT}). El bloque
\texttt{recommendation\_funnel} del trace expone el embudo. \\

Reranking estructurado por \texttt{providerFitCriteria} &
El extractor emite criterios de ajuste (presupuesto, tipo de evento,
preferencias, restricciones duras) que el runtime aplica de forma
determinista sobre el detalle del proveedor. El agente conversacional
no rerankea proveedores: sólo decide cómo presentar el shortlist ya
ordenado. \\

Manejo de múltiples selecciones en un mismo turno &
Cuando un usuario combina ``quiero al DJ Pulga y ahora muéstrame
catering'', la extracción incluye \texttt{secondaryIntents} y la
selección del DJ se preserva aunque la intención principal sea
\texttt{buscar\_proveedores}. La pista de selección se autocompleta
por heurística de alias cuando el extractor no la emite. \\

Nodos informativos de primera clase &
\texttt{consultar\_faq} y \texttt{consultar\_evento\_invitado} son
nodos del grafo, no herramientas sueltas. Tienen su propio bundle de
prompts, herramientas dedicadas (\texttt{file\_search} en
\texttt{consultar\_faq}, \texttt{lookup\_user\_event\_context} en
\texttt{consultar\_evento\_invitado}) y rutas claras de entrada y
salida. FAQ exige la invocación real del tool antes de responder. \\

Cierre con cotización real (\texttt{finish\_plan}) &
La herramienta \texttt{finish\_plan} itera sobre todas las
necesidades con proveedor seleccionado y llama a
\texttt{POST /api-web/vendor/quote} por cada uno. Devuelve un
\texttt{status} agregado (\texttt{success}, \texttt{partial},
\texttt{failed}) y, en éxito, muta el plan a
\texttt{lifecycle\_state: finished} y a
\texttt{current\_node: necesidad\_cubierta}. \\

Telemetría por turno siempre activa &
Cada turno produce un \texttt{TurnTrace} con tiempos por etapa,
herramientas, parámetros, errores, embudo de recomendación y
resúmenes estructurados (decisión, extracción, plan). El registro
se persiste en la tabla de perf con TTL configurable y se devuelve
al cliente sólo cuando \texttt{client\_mode=cli}. \\

Renderers por canal deterministas &
La salida del agente se parsea en estructuras tipadas
(\texttt{welcome}, \texttt{recommendation}, etc.) y se proyecta a
WhatsApp (\texttt{*negrita*} y emojis), webchat (texto plano con
viñetas y URLs) o terminal. El formato nunca queda a criterio del
modelo. \\

CLI de desarrollo que emula WhatsApp &
La terminal apunta siempre a la Function URL desplegada (no corre un
agente local), emula el contrato del canal, expone trace, plan y
perf en cada turno, soporta \texttt{/help}, \texttt{/config},
\texttt{/plan} y \texttt{/exit}, y muestra una línea de progreso
\textit{in-place} con fases. \\

\bottomrule
\caption{Resumen de características clave.}
\end{longtable}

% ============================================================
\section{Decisiones de diseño relevantes}
% ============================================================

\subsection{Determinismo de ruteo y parseo canónico}

El proyecto se aleja explícitamente del \emph{matching} de cadenas o
palabras clave en favor de evidencia estructurada validada por Zod. El
log documenta múltiples decisiones consecutivas en este sentido:

\begin{itemize}[leftmargin=*,itemsep=1pt]
  \item Normalización canónica de tipos de evento, niveles de precio,
        categorías de proveedor, claves de país y nodos de decisión.
  \item Eliminación de la autoridad del modelo sobre
        \texttt{actions} generadas: el control de flujo queda en
        manos de la combinación de \texttt{intent}, pista de selección,
        estado del nodo y estado persistente.
  \item \texttt{providerFitCriteria} como contrato estructurado para
        reordenar proveedores: el extractor lo emite y el runtime lo
        aplica de forma determinista.
  \item Esquemas Zod para \emph{close actions}, \emph{provider
        references}, \emph{plan operations}, \emph{query intents},
        \emph{explanations} y \emph{detail requests} en
        \texttt{src/runtime/extraction-schemas.ts}.
\end{itemize}

\subsection{\emph{TurnDecision} como contrato de ruteo}

La conversión del grafo implícito en un objeto
\texttt{TurnDecision} validado por Zod resolvió un bug persistente: el
\texttt{active\_need} quedaba obsoleto tras turnos multi-frente, lo que
provocaba que el sistema ignorase nuevas necesidades para centrarse
sólo en la más antigua. Con \texttt{TurnDecision}, el modelo interpreta
y el código enruta; los tests pueden asegurar que
\texttt{turn\_decision.nextNode} coincide con el nodo ejecutado.

\subsection{Persistencia \emph{event-plan-first} con migración limpia}

El plan migró de campos singulares (\texttt{selected\_provider\_id},
\texttt{selected\_provider\_hint}) a arreglos. \texttt{normalizeRawPlan}
realiza la migración en el \emph{boundary} de carga, lo que evita
mantener un shim de retrocompatibilidad en el código de aplicación y
permite que la purga de planes antiguos sea un paso explícito.

\subsection{Embedding, prompt cache y retries}

El runtime apunta a la familia \texttt{gpt-5.4} con dos modelos
(\texttt{gpt-5.4-mini} para respuesta y \texttt{gpt-5.4-nano} para
extracción) y configura:

\begin{itemize}[leftmargin=*,itemsep=1pt]
  \item \texttt{promptCacheRetention} (\texttt{in-memory} o
        \texttt{24h}) con un \texttt{cacheKey} estable por bundle.
  \item \texttt{maxRetries: 3} en el cliente OpenAI y
        \texttt{retryPolicies.any(retryPolicies.httpStatus([429]),
        retryPolicies.networkError())} con backoff exponencial y
        jitter ($1\,\mathrm{s} \to 30\,\mathrm{s}$) para tolerar ventanas
        de TPM de 20 segundos.
  \item Snapshots compactos del plan en extractores y composición de
        respuesta para reducir tokens.
  \item \texttt{reasoning.effort: none} y \texttt{text.verbosity: low}
        en el runtime de respuesta para minimizar latencia.
\end{itemize}

\subsection{Canales agnósticos y \texttt{client\_mode}}

El núcleo sólo conoce \texttt{channel} y \texttt{user\_id}; no contiene
lógica de WhatsApp. La respuesta se serializa siempre con un campo
\texttt{message} compatible con cualquier cliente. Los diagnósticos
(\texttt{trace}, \texttt{perf}, \texttt{plan}) se devuelven únicamente
cuando el cliente declara \texttt{client\_mode: "cli"}, de modo que
los canales de consumo no ven datos internos. Los adaptadores son
responsables de la autenticación de webhook, la idempotencia, el
formateo final y la captura de feedback.

\subsection{Telemetría de bajo costo}

Las trazas se persisten en una tabla DynamoDB con TTL
(\texttt{ttl\_epoch\_seconds}) y un GSI para consultas por canal/usuario.
Los registros contienen resúmenes estructurados (\texttt{extraction\_summary},
\texttt{plan\_summary}, \texttt{recommendation\_funnel}) y un
\texttt{user\_message\_preview} truncado, suficientes para reconstruir
interacciones fallidas sin almacenar blobs crudos.

\subsection{Prompts versionados y migraciones limpias}

El proyecto no mantiene shims de retrocompatibilidad: cada cambio
sustancial va acompañado de una migración atómica en planes, fixtures
y casos de evaluación. Los prompts viven como texto en español en
\texttt{prompts/} y el runtime los carga por nodo, lo que permite
auditar la trazabilidad de cada cambio en el log de implementación.

% ============================================================
\section{Sistema de prompts y conocimiento}
% ============================================================

\subsection{Estructura de \texttt{prompts/}}

\begin{itemize}[leftmargin=*,itemsep=1pt]
  \item \texttt{shared/}: sistema base, alcance, conocimiento de
        dominio, estilo, disciplina de flujo, estrategia de preguntas,
        anti-patrones comunes.
  \item \texttt{nodes/<nodo>/system.txt}, \texttt{response\_contract.txt},
        \texttt{tool\_policy.txt}: contrato por nodo.
  \item \texttt{extractors/}: bundle específico para el agente extractor
        (\texttt{field\_definitions}, \texttt{normalization\_rules},
        \texttt{conflict\_resolution}, \texttt{examples}, etc.).
\end{itemize}

La carga se hace a través de \texttt{PromptLoader}
(\texttt{src/runtime/prompt-loader.ts}) y se cachea por bundle para
aprovechar la caché de prompts de OpenAI.

\subsection{Capa de conocimiento (FAQ)}

El sistema integra una \emph{base de conocimiento} derivada de la
\textit{help center} pública de Sin Envolturas. El proceso completo
está descrito en \texttt{docs/knowledge-base-integration.md}:

\begin{enumerate}[leftmargin=*,itemsep=1pt]
  \item \texttt{TawkHelpScraper} descarga los artículos de
        \texttt{https://sinenvolturas.tawk.help}.
  \item El formateador produce un archivo Markdown por artículo con
        \emph{frontmatter} YAML (\texttt{title}, \texttt{slug},
        \texttt{category}, \texttt{article\_type}, \texttt{tags},
        \texttt{source\_url}, \texttt{last\_updated},
        \texttt{related\_topics}).
  \item \texttt{OpenAiKnowledgeUploader} sube los artículos a un vector
        store de OpenAI en lotes (\texttt{uploadBatch}) y elimina los
        archivos de lotes anteriores (\texttt{cleanupOldBatches}).
  \item Un Lambda de sincronización (\texttt{recap-agent-knowledge-sync-dev})
        corre semanalmente con \texttt{rate(7 days)} desde
        EventBridge.
  \item El nodo \texttt{consultar\_faq} recibe la herramienta
        hospedada \texttt{file\_search} restringida a ese vector
        store.
\end{enumerate}

En la práctica, la descarga inicial tuvo que hacerse localmente porque
Tawk bloquea las IP de AWS Lambda y GitHub Actions.

% ============================================================
\section{Marco de evaluación}
% ============================================================

El sistema incluye un \emph{marco de evaluación nativo} descrito en
\texttt{docs/evaluation-framework.md} y \texttt{src/evals/}. Sus
principios son:

\begin{itemize}[leftmargin=*,itemsep=1pt]
  \item \textbf{Casos versionados en YAML}: \texttt{evals/cases/}
        contiene escenarios; \texttt{evals/templates/} expone bases
        reutilizables; \texttt{evals/fixtures/} tiene planes semilla
        y fragmentos de respuesta de proveedores.
  \item \textbf{Dos targets}: \texttt{offline} (in-memory,
        determinista, ideal para CI) y \texttt{live\_lambda} (invoca
        la Function URL desplegada y normaliza la respuesta al mismo
        sobre de resultado).
  \item \textbf{Expectativas por capas}: campos de plan, transiciones
        de nodo, herramientas invocadas, conteo de proveedores, campos
        de traza, comparaciones tolerantes de texto y jueces
        semánticos opcionales.
  \item \textbf{Métricas SOTA}: \texttt{tool\_precision},
        \texttt{tool\_recall}, \texttt{tool\_F1}, cobertura de ramas,
        tasas de paso de estado y trayectoria, tasa de persistencia,
        tasa de acierto de caché, distribución de latencia y totales
        de tokens.
  \item \textbf{Artefactos para BI}: cada corrida emite
        \texttt{report.json}/\texttt{report.md},
        \texttt{dashboard.json} y \texttt{dashboard.csv} bajo
        \texttt{.eval-runs/<run-id>/}.
  \item \textbf{Concurrencia}: el runner soporta
        \texttt{--parallel <n>} con un pool acotado que preserva el
        orden determinista de los resultados.
  \item \textbf{Eval observable}: \texttt{eval:observable-live} es un
        runner que imprime sólo la transcripción (turno del usuario y
        respuesta del agente) usando un planificador con estado y
        turnos barajados, para inspección manual de la conversación
        real.
\end{itemize}

Los \emph{suites} principales son \texttt{smoke}, \texttt{dev\_regression},
\texttt{benchmark\_full}, \texttt{feedback\_regression},
\texttt{live\_smoke}, \texttt{live\_comprehensive} y
\texttt{live\_feedback\_token\_regression}.

% ============================================================
\section{Evolución del proyecto}
% ============================================================

El log de implementación (\texttt{docs/implementation-log.md}) recoge
cronológicamente los hitos del proyecto. La tabla~\ref{tab:hitos}
resume los más representativos.

\begin{longtable}{p{2.4cm}p{4.0cm}p{9.0cm}}
\toprule
\textbf{Fecha} & \textbf{Hito} & \textbf{Resumen} \\
\midrule
\endhead

2026-04-05 & Esqueleto del runtime &
Convenciones en \texttt{AGENTS.md}, scaffolding de TypeScript,
arquitectura basada en grafo de nodos, \texttt{PlanStore} con
DynamoDB e in-memory, primer SDK de OpenAI Agents y gateway real
contra Sin Envolturas. \\

2026-04-05 & Secreto en AWS Secrets Manager &
La API key de OpenAI se resuelve en tiempo de ejecución desde
Secrets Manager; \texttt{deploy.mjs} sincroniza el secreto desde
\texttt{.env}. \\

2026-04-06 & Prompts por nodo &
Los prompts pasan de ser texto suelto a bundles estructurados
(\texttt{system}, \texttt{response\_contract}, \texttt{tool\_policy}),
compartidos entre conversacional y extractor, con alcance de
herramientas por nodo. \\

2026-04-07 & Configuración centralizada y prohibición de \texttt{any} &
\texttt{src/runtime/config.ts} se vuelve la fuente de verdad para
modelos, límites de búsqueda y canal por defecto; ESLint y TS
aplican el veto de \texttt{any} explícito. \\

2026-04-07 & \texttt{nodejs24.x} de extremo a extremo &
Local, Lambda y \texttt{.nvmrc} se alinean en Node 24 LTS. \\

2026-04-08 & \emph{Event-plan-first} &
El plan admite múltiples \texttt{provider\_needs} con
\texttt{active\_need\_category}. Se reescriben extractor y prompts
para razonar sobre el evento y no sobre un proveedor puntual. \\

2026-04-12 & Marco de evaluación nativo &
Aparecen \texttt{src/evals/} y los activos en \texttt{evals/}, con
expectativas en capas y dos targets (offline y live). \\

2026-04-14 & Búsqueda resiliente a ubicación dispersa &
El gateway prefiere coincidencias exactas de ciudad, pero recurre a
candidatos de la categoría con metadatos más amplios cuando la
granularidad de la ubicación es insuficiente. \\

2026-04-20 & Búsqueda mixta de proveedores y embudo 15 $\to$ 5 &
Se consultan \texttt{/filtered} y \texttt{/filtered/full}, se
fusionan los resultados y se entrega al modelo de respuesta un
top-15, restringido a un top-5 visible. \\

2026-04-20 & \texttt{finish\_plan}, ciclo de vida y TTL &
Aparecen \texttt{lifecycle\_state}, \texttt{contact\_name},
\texttt{contact\_email} y la herramienta
\texttt{finish\_plan}, con TTL de 24\,h en el plan
\texttt{finished} (luego retirado). \\

2026-04-23 & Cotización real y separación pausa/cierre &
\texttt{finish\_plan} llama a
\texttt{POST /api-web/vendor/quote} por proveedor seleccionado;
\texttt{pausar} se separa de \texttt{cerrar} y el nodo
\texttt{crear\_lead\_cerrar} ejecuta un flujo de tres pasos
(contacto $\to$ resumen $\to$ confirmación). \\

2026-04-28 & Integración real de la KB de Tawk &
Se conecta el vector store con \texttt{file\_search}, se rediseña
el scraper (un archivo por artículo con metadatos), se crea el
stack \texttt{recap-agent-knowledge-sync-dev} y se documenta la
limitación de IP. \\

2026-04-30 & Renderers por canal estructurados &
Las respuestas se parsean en estructuras tipadas antes de
renderizarse en WhatsApp o webchat. La terminal emula WhatsApp. \\

2026-05-05 & Heurísticas de selección multi-intención &
Aparecen \texttt{secondaryIntents},
\texttt{resolveEffectiveSelectionHint} y el prompt
correspondiente para evitar la pérdida de
\texttt{selectedProviderHint} en turnos que combinan selección con
apertura de nueva necesidad. \\

2026-05-05 & Búsqueda vectorial de proveedores &
Se añade \texttt{provider-sync} (scrape + formato + upload) y el
modo \texttt{hybrid} en el gateway. \\

2026-05-06 & Categorías canónicas y filtros vectoriales &
Las categorías dejan de tener alias heurísticos y pasan a ser un
enum Zod; los filtros del vector store usan la clave canónica
exacta, lo que recupera a proveedores como Dulcefina. \\

2026-05-07 & Múltiples proveedores por necesidad &
Los campos \texttt{selected\_provider\_*} pasan a arreglos; la
resolución de selección admite ordinales, alias y descriptores.
\texttt{finish\_plan} crea cotizaciones por cada proveedor
seleccionado. \\

2026-05-14 & Normalización canónica final &
Tipos de evento, niveles de precio, claves de país y nodos de
decisión se consolidan en módulos canónicos. Se eliminan las
\texttt{actions} generadas como ruta de control. \\

2026-05-22 & Nodo \texttt{consultar\_evento\_invitado} y telemetría ampliada &
Se añade el modo de consulta de eventos como invitado
(\texttt{lookup\_user\_event\_context}) y se amplían los
\texttt{TurnTrace} y \texttt{TurnPerfRecord} con resúmenes
estructurados. \\

2026-06-04 & \texttt{TurnDecision} autoritativo y búsqueda por sub-queries &
La decisión del turno pasa a ser el contrato principal de ruteo; se
introduce la búsqueda por sub-consulta dentro de una misma
necesidad para no colapsar ``sushi + torta'' en un solo shortlist. \\

\bottomrule
\caption{Hitos del proyecto.}
\label{tab:hitos}
\end{longtable}

% ============================================================
\section{Entorno, despliegue y operaciones}
% ============================================================

\subsection{Entorno local y comandos}

El repositorio requiere \texttt{node$\geq$24} y Bun opcional. Los
comandos principales son:

\begin{codeblock}
npm install
npm run typecheck   # tsc --noEmit
npm run lint        # eslint
npm test            # vitest run
npm run check       # typecheck + lint + test
npm run build       # esbuild -> dist/
npm run deploy      # scripts/deploy.mjs (Lambda + CloudFormation)
npm run terminal    # CLI contra la Function URL desplegada
npm run eval        # suite offline/live
npm run eval:observable-live  # transcript barajado en vivo
\end{codeblock}

\noindent\textbf{Cuadro:} Comandos principales del repositorio.

\subsection{Configuración tipada}

\texttt{src/runtime/config.ts} define un esquema Zod sobre variables
de entorno y un objeto \texttt{AppConfig} tipado que el handler y los
componentes consumen. La configuración cubre:

\begin{itemize}[leftmargin=*,itemsep=1pt]
  \item OpenAI: API key, \texttt{OPENAI\_SECRET\_ID}, modelos,
        \texttt{promptCacheRetention}.
  \item AWS: región, tablas DynamoDB.
  \item Sin Envolturas: \texttt{SINENVOLTURAS\_BASE\_URL},
        \texttt{SINENVOLTURAS\_GUEST\_SERVICE\_BASE\_URL}, modo de
        búsqueda, vector store.
  \item Recomendación: \texttt{REPLY\_PROVIDER\_LIMIT},
        \texttt{PRESENTATION\_PROVIDER\_LIMIT},
        \texttt{PROVIDER\_DETAIL\_LOOKUP\_LIMIT}.
  \item Telemetría: \texttt{PERF\_TABLE\_NAME},
        \texttt{PERF\_RETENTION\_DAYS}.
  \item Knowledge base: \texttt{KB\_ENABLED},
        \texttt{KB\_BASE\_URL}, \texttt{KB\_VECTOR\_STORE\_ID}.
  \item \emph{Feature flags} del agente:
        \texttt{AGENT\_FEATURE\_PROVIDER\_PLANNING},
        \texttt{\_SEARCH}, \texttt{\_QUOTE\_REQUESTS},
        \texttt{\_FAQ}, \texttt{\_INVITED\_EVENT\_LOOKUP}.
\end{itemize}

Los \emph{feature flags} controlan, por ejemplo, el menú de bienvenida
dinámico, lo que permite activar o desactivar capacidades sin tocar
prompts.

\subsection{Despliegue}

\begin{enumerate}[leftmargin=*,itemsep=1pt]
  \item \texttt{node scripts/build.mjs} produce \texttt{dist/}
        con \texttt{lambda/index.js},
        \texttt{knowledge-sync/index.js}, etc.
  \item \texttt{scripts/deploy.mjs} sube el artefacto a S3 y
        despliega la plantilla de CloudFormation.
  \item El secreto de OpenAI se sincroniza desde \texttt{.env} a
        Secrets Manager; el Lambda lo lee en tiempo de ejecución.
\end{enumerate}

El endpoint público de desarrollo documentado es

\begin{center}
\texttt{https://jwtjjociscvaa5dsrp5gokmno40doiva.lambda-url.us-east-1.on.aws/}
\end{center}

con stack \texttt{recap-agent-runtime}, tablas
\texttt{recap-agent-runtime-plans} y
\texttt{recap-agent-runtime-perf}.

\subsection{Operación}

\begin{itemize}[leftmargin=*,itemsep=1pt]
  \item Purga de planes de terminal: \texttt{npm run purge:terminal-plans}
        con \texttt{--dry-run} o \texttt{--yes}.
  \item Sync manual de proveedores: \texttt{npm run sync:providers}.
  \item Sync manual de la base de conocimiento:
        \texttt{KB\_SKIP\_UPLOAD=true npx tsx scripts/sync-knowledge-base.ts}.
\end{itemize}

% ============================================================
\section{Limitaciones conocidas y trabajo futuro}
% ============================================================

\begin{itemize}[leftmargin=*,itemsep=2pt]
  \item \textbf{Tawk bloquea AWS Lambda y GitHub Actions}, por lo que
        la descarga inicial de la base de conocimiento debe hacerse
        desde una máquina local; el \emph{Lambda} sólo se encarga del
        ciclo de subida al vector store.
  \item El scraping y la sincronización inicial del vector store de
        proveedores requieren intervención manual para crear el
        recurso en OpenAI; existe un comando
        \texttt{npm run sync:providers} para poblarlo después.
  \item \textbf{Streaming fuera de alcance}: el canal de WhatsApp no
        soporta respuestas en flujo. Por convención del proyecto, el
        terminal emula WhatsApp y no introduce capacidades nuevas.
  \item El webhook de WhatsApp y la autenticación de los adaptadores
        se mantienen intencionalmente fuera de \texttt{src/runtime/};
        son responsabilidad de cada adaptador de canal.
  \item El catálogo canónico de categorías y tipos de evento requiere
      reindexar el vector store y los planes antiguos cuando se
      amplían (cambios disruptivos sin \emph{shims} de
      retrocompatibilidad).
\end{itemize}

% ============================================================
\section{Conclusiones}
% ============================================================

\texttt{recap-agent} demuestra que es posible construir un agente
conversacional serverless, multicanal y de baja operación, sin
sacrificar la trazabilidad ni la calidad de las decisiones de flujo.
Los principios que el proyecto ha consolidado a lo largo de su
evolución son:

\begin{itemize}[leftmargin=*,itemsep=1pt]
  \item \textbf{El modelo interpreta, el código decide}: extracción
        estructurada + \texttt{TurnDecision} determinista.
  \item \textbf{Plan de evento primero, no proveedor}: las
        necesidades múltiples, los sub-queries y las prioridades por
        tipo de evento se modelan como ciudadanos de primera clase.
  \item \textbf{Prompts como contrato}: cada nodo define su
        comportamiento vía bundles versionados en español.
  \item \textbf{Canales y almacenamiento como adaptadores delgados}:
        el núcleo permanece agnóstico y testeable.
  \item \textbf{Evaluación por capas}: casos reutilizables,
        expectativas estructuradas y métricas SOTA en el mismo marco.
  \item \textbf{Telemetría barata y útil}: trazas por turno con TTL
        y resúmenes estructurados, expuestas al cliente sólo en modo
        \texttt{cli}.
\end{itemize}

El resultado es un sistema que puede desplegarse con
\texttt{npm run deploy}, ejercitarse con \texttt{npm run terminal} o
con \texttt{npm run eval}, y extenderse añadiendo un nuevo nodo al
grafo, un nuevo \emph{feature flag} o un nuevo canal, sin reescribir
el núcleo.

% ============================================================
\appendix
\section{Apéndice A: Variables de entorno clave}
% ============================================================

\begin{longtable}{p{6.0cm}p{3.5cm}p{6.0cm}}
\toprule
\textbf{Variable} & \textbf{Por defecto} & \textbf{Propósito} \\
\midrule
\endhead

\texttt{OPENAI\_MODEL} & \texttt{gpt-5.4-mini} & Modelo de respuesta. \\
\texttt{OPENAI\_EXTRACTOR\_MODEL} & \texttt{gpt-5.4-nano} & Modelo extractor. \\
\texttt{OPENAI\_PROMPT\_CACHE\_RETENTION} & \texttt{in-memory} & Retención de caché de prompts. \\
\texttt{AWS\_REGION} & \texttt{us-east-1} & Región AWS. \\
\texttt{PLANS\_TABLE\_NAME} & \texttt{recap-agent-plans} & Tabla DynamoDB de planes. \\
\texttt{PERF\_TABLE\_NAME} & \texttt{recap-agent-runtime-perf} & Tabla de telemetría. \\
\texttt{PERF\_RETENTION\_DAYS} & \texttt{30} & TTL de la tabla de perf. \\
\texttt{PROMPTS\_DIR} & \texttt{/var/task/prompts} & Ruta al bundle de prompts. \\
\texttt{SINENVOLTURAS\_BASE\_URL} &
  \texttt{https://api.sinenvolturas.com/api-web/vendor} &
  Base de la API del marketplace. \\
\texttt{SINENVOLTURAS\_GUEST\_SERVICE\_BASE\_URL} &
  \texttt{https://api.sinenvolturas.com/api/guest-service} &
  Base del servicio de invitados. \\
\texttt{DEFAULT\_INBOUND\_CHANNEL} & \texttt{terminal\_whatsapp} & Canal por defecto. \\
\texttt{PROVIDER\_SEARCH\_LIMIT} & \texttt{12} & Candidatos persistidos. \\
\texttt{SEARCH\_SUMMARY\_WORD\_LIMIT} & \texttt{5} & Resumen de búsqueda. \\
\texttt{REPLY\_PROVIDER\_LIMIT} & \texttt{6} & Top que recibe el LLM. \\
\texttt{PRESENTATION\_PROVIDER\_LIMIT} & \texttt{6} & Top visible al usuario. \\
\texttt{PROVIDER\_DETAIL\_LOOKUP\_LIMIT} & \texttt{3} & Detalles por turno. \\
\texttt{PROVIDER\_SEARCH\_MODE} & \texttt{hybrid} & \texttt{api}/\texttt{vector}/\texttt{hybrid}. \\
\texttt{PROVIDER\_VECTOR\_STORE\_ID} & (vacío) & ID del vector store de proveedores. \\
\texttt{KB\_ENABLED} & \texttt{true} & Habilita la KB. \\
\texttt{KB\_VECTOR\_STORE\_ID} & (vacío) & ID del vector store de la KB. \\
\texttt{AGENT\_FEATURE\_*} & \texttt{true} & Feature flags por capacidad. \\

\bottomrule
\caption{Principales variables de entorno del runtime.}
\end{longtable}

% ============================================================
\section{Apéndice B: Estructura del repositorio}
% ============================================================

\begin{codeblock}
recap-agent/
  AGENTS.md
  README.md
  docs/
    implementation-log.md
    channel-integration.md
    evaluation-framework.md
    knowledge-base-integration.md
    feedback-implementation-plan.md
    feedback-test-coverage.md
    provider-vector-search.md
  infra/
    cloudformation/stack.yaml
    knowledge-sync.yml
    provider-sync.yml
  prompts/
    shared/         (base_system, output_style, ...)
    nodes/<nodo>/   (system, response_contract, tool_policy)
    extractors/     (system, field_definitions, ...)
  src/
    core/           (decision-nodes, plan, sufficiency, turn-decision, ...)
    runtime/        (agent-service, openai-agent-runtime, gateways, ...)
    storage/        (plan-store, perf-store, adapters Dynamo/in-memory)
    lambda/handler.ts
    terminal/client.ts
    knowledge-sync/ (scraper, formatter, uploader, sync, handler)
    provider-sync/  (fetcher, formatter, uploader, sync, handler)
    evals/          (case-schema, runner, reporting, targets, scorers)
    logs/trace/perf.ts
  evals/
    cases/          (escenarios YAML)
    templates/      (plantillas base)
    fixtures/       (planes semilla y respuestas de proveedores)
    suites/         (smoke, dev_regression, benchmark_full, ...)
    matrices/       (matrices de modelos y configuraciones)
  scripts/          (build, deploy, purge, sync-*)
  tests/            (vitest)
  package.json
  tsconfig.json
  eslint.config.mjs
  .env.example
\end{codeblock}

\noindent\textbf{Cuadro:} Estructura principal del repositorio.

\end{document}
